% V18 integration of Dot batch 090; immutable source SHA256 e80656b5f3b296d20bf542ed36b5e58f6cd149b998dc84e1a8dc17517777fcbf
\begingroup
% These definitions have only group-local scope; the shared preamble is untouched.
\def\OO{\mathcal O}\def\FF{\mathcal F}\def\GG{\mathcal G}
\def\PP{\mathbf P}\def\calP{\mathcal P}\def\dd{\delta}
\def\Div{\operatorname{Div}}\def\depth{\operatorname{depth}}
\def\pr{\operatorname{pr}}\def\coker{\operatorname{coker}}
记 $P=\PP^n_k$，$P^\vee$ 为超平面参数空间，$\eta$ 为其泛点，
$K_\eta=k(\eta)$，$h:Y_\eta\to X$ 为泛超平面截面的投影。


\egas{2.14}{定理}
在（2.12）的非分歧情形，设 $y\in Y_\eta$，$x=h(y)$。
对域 $K$ 上的诺特局部代数 $A$，令 $\mathcal P(A,K)$ 为下列任意一种性质：
\begin{enumerate}
 \item[(i)] $A$ 在 $K$ 上几何正则；
 \item[(ii)] $A$ 在 $K$ 上几何满足 $(R_r)$；
 \item[(iii)] $A$ 在 $K$ 上可分；
 \item[(iv)] $A$ 在 $K$ 上几何正规。
\end{enumerate}
于是有下列等价\footnote{本文用 $r$ 表示 Serre 条件的整数下标，以区别底域 $k$；法、英稿写 $(R_k)$。
英稿将右端底域排作 $k(y)$；法本手写符号与 $Y$ 的泛参数下标及其后 $k(\eta)$ 的写法一致，
也和泛纤维自然结构相符。$\OO_{Y_\eta,y}$ 有到剩余域 $k(y)$ 的商映射，
而不是由此获得自然的 $k(y)$-代数结构。法本左边批语受扫描阴影遮挡，
英编稿注8读作“非分歧，或特征零”，且主语录作 $X$ 并带编辑标记。
这里把非分歧赋予 $f$；特征零的局部结论及其边界在后文单列，未以边批扩大全部推论的范围。}：
\[
 \mathcal P(\OO_{X,x},k)\ \Longleftrightarrow\
 \mathcal P(\OO_{Y_\eta,y},K_\eta).
\]
\begin{proof}[原稿论证]
原稿先说明：（ii）和（2.8）（ii）结合第6段的结果，给出（iii）、（iv）；
（i）已由（2.12）证明。由（i）推（ii），令 $Z$ 为 $X$ 对 $k$ 的非光滑轨迹。
由（2.12），$h^{-1}(Z)$ 恰是 $Y_\eta$ 对 $K_\eta$ 的非光滑轨迹。
原稿援引平坦性说明，两坏轨迹在相应点的局部余维相等，故一边大于 $r$ 当且仅当另一边大于 $r$。
\end{proof}

\noindent\emph{译校补证。}
这里局部余维是局部谱中坏闭集各泛点的最小高度，空集取 $+\infty$；
它不是任意一对对应点的高度。设 $A\to B$ 为诺特平坦局部同态，$I\subseteq A$ 定义该坏闭集。
任取 $\mathfrak q\supseteq IB$，其收缩 $\mathfrak p$ 含 $I$，下降定理给出
$\operatorname{ht}\mathfrak q\ge\operatorname{ht}\mathfrak p$。
反之，选取 $I$ 上高度最小的素理想 $\mathfrak p$，由忠实平坦性选到其上方的素理想，
再在其中取 $\mathfrak pB$ 上的极小素 $\mathfrak q$。
其收缩仍为 $\mathfrak p$，且局部纤维 $B_{\mathfrak q}/\mathfrak p B_{\mathfrak q}$ 零维，
依高度公式有 $\operatorname{ht}\mathfrak q=\operatorname{ht}\mathfrak p$。
若 $I=A$，两坏集同时为空。这就给出所需的最小高度等式。

有限型域上概形的几何 $(R_r)$ 可在底域代数闭包上检查。
在任一提升点，非光滑轨迹为原坏集的逆像，又因代数闭域完美，该坏集恰是非正则轨迹。
对各提升点的平坦局部基变换用上述最小高度等式，故“在高度不超过 $r$ 的素点都正则”
等价于原局部谱中非光滑轨迹的最小高度大于 $r$。
由此得到（ii），而非将原底域上的普通正则性误作几何正则性。
环的“可分”在此读作几何约化；几何约化与几何正规分别使用
$(R_0)+(S_1)$ 与 $(R_1)+(S_2)$。代数扩域的局部纤维零维，
诺特平坦局部映射的深度、维数公式给出这些 $(S_i)$ 条件在该类扩域前后一致；
再用（2.8）（ii）比较原局部谱的 $(S_i)$。这段补证限于这里的结构环，不代替（2.16）的模层论证。
\footnote{法稿两处均援引第6段；英编稿注9指 EGA IV 的 6.5，属编辑线索，未据此倒填原稿卷次。
现代表述见 \href{https://stacks.math.columbia.edu/tag/00ON}{Stacks 00ON}（高度公式）与
\href{https://stacks.math.columbia.edu/tag/033P}{033P}（Serre 条件）。}

\egas{2.15}{推论}
仍设 $f$ 非分歧。设 $Z$ 为 $X$ 不在 $k$ 上光滑（分别地，
不在 $k$ 上几何满足 $(R_r)$、不可分、不几何正规）的点集。
则 $Y_\eta$ 在 $K_\eta$ 上光滑（分别地，几何满足 $(R_r)$、可分、几何正规），
当且仅当 $Z$ 有限。\footnote{法13在“正规”之前清楚地手补了“géom.”；英稿左侧漏译，
右侧则有“几何正规”。推论用到（2.3）和有限型非分歧态射的拟有限性：
泛截面避开坏集首先等价于 $f(Z)$ 有限，继而才等价于 $Z$ 有限。
这里只在坏轨迹构造性成立的范围使用（2.3）；对 $(R_r)$，
可由坏闭集的有限个不可约分支中高度不超过 $r$ 的那些取闭包并集得到相应坏轨迹。}

从行文上看，（2.14）、（2.15）已经包含（2.12）、（2.13），后两条可以不必另行陈述。
另一方面，推论在应用中更重要；可将定理放作预备命题或引理，突出推论。
对性质（iii），还可作如下变体。

\egas{2.16}{推论}
仍设 $f$ 是浸入，而且 $\FF$ 凝聚。在（2.7）的条件下，
$Z_i$ 不是嵌入关联分支，当且仅当相应的 $Z_{i\eta}$ 不是嵌入关联分支。
在这种情况下，$\FF$ 在 $Z_i$ 处相对于 $k$ 的根重数，
等于 $\GG_\eta$ 在对应分支处相对于 $K_\eta$ 的根重数（原稿参考待补）。
\footnote{法本 \emph{non immergé} 在这里修饰关联分支，不能译作“不是浸入态射”。
第二个分支下标在法13正文不完整；由（2.7）两边的对象及英稿第一句恢复为对应泛截面分支。
“根重数”保留旧稿 \emph{multiplicité radicielle} 的称呼，不把它直接替换为普通局部长度。}
\begin{proof}[原稿论证]
第一断言见（2.7）的最后一条。对第二断言，由 $h$ 平坦，函子 $\FF\mapsto\GG_\eta$ 正合。
限制到 $Z_i$ 泛点的一个邻域并逐层分解，归约到 $\FF=\OO_{Z_i}$，可设 $Z_i=X$。
也可以先归约到 $k$ 代数闭；此时 $X$ 在 $k$ 上可分。然后，原稿说该结论包含在（2.15）的（iii）中。
像通常那样，可推出下一条。
\end{proof}
\noindent\emph{译注。} 法13明写 \emph{séparable}（可分），非英稿的 \emph{separate}（分离）；
手稿括弧涉及代数闭归约和此后的可分性，两句均应保留。
原稿没有在此给出根重数的定义引用或交代扩域前后重数的归约，故未把这段补写为独立证明。

\egas{2.17}{推论}
仍沿用前述模层设定。令 $Z$ 为 $X$ 上 $\FF$ 不对 $k$ 可分的点集（参考待补）。
则 $\GG_\eta$ 对 $K_\eta$ 可分，当且仅当 $Z$ 有限。
特别地，若 $\FF$ 对 $k$ 可分，则 $\GG_\eta$ 对 $K_\eta$ 可分。

\egas{2.18}{备注}
若完全去掉对 $f$ 的限制，关键结果（2.12）及由此所得（2.13）、（2.17）
一般不再成立。例如可让 $X$ 为 $P$ 的纯不可分覆盖。
当底域特征为零时，局部光滑的等价（2.12）不需假设 $f$ 是浸入。
\footnote{法14此处写“2.12 至 2.17”，末位有改写痕迹；英稿取“2.12 和 2.17”。
原稿继称只需核对2.12，并引2.11以及特征零下光滑等于正则。
这个简写不能作为把所有涉及坏集有限性的命题推广到任意 $f$ 的证明；
模层（2.17）的特征零推广须另按其定义和假设核对。}
特征零时，$k$ 与 $K_\eta$ 皆完美，故两边有限型概形的光滑性可按局部环正则性检查；
$h$ 是正则态射（2.8），因此正则性沿该平坦局部同态下降，并凭正则纤维上升。

\par\noindent\textit{译校补证另见核查报告。}\cmt{D18-090-R}{纯不可分覆盖及特征零下常值态射的两个反例完整保存在本版核查报告“D18-090”条；其中译校补证为非原文。}
\noindent 法14在2.18后另有划销的过渡句，意为从已经处理的局部性质转向整体性质；
不列入正文。第3节及其页下方划销草稿的校记接在下文。

\endgroup
